% Generated draft: requires mathematical and source review.
\documentclass[11pt]{article}
\usepackage[margin=25mm]{geometry}
\usepackage{fontspec}
\setmainfont{Latin Modern Roman}
\usepackage{amsmath,amssymb,mathtools,unicode-math}
\setmathfont{Latin Modern Math}
\usepackage{enumitem,needspace,xurl,hyperref}
\setlength{\parindent}{0pt}
\setlength{\parskip}{5pt}
\setlength{\emergencystretch}{3em}
\newcommand{\N}{\mathbb{N}}
\newcommand{\Z}{\mathbb{Z}}
\newcommand{\Q}{\mathbb{Q}}
\newcommand{\R}{\mathbb{R}}
\newcommand{\C}{\mathbb{C}}
\newcommand{\F}{\mathbb{F}}
\newcommand{\A}{\mathbb{A}}
\newcommand{\PP}{\mathbb{P}}
\newcommand{\E}{\mathbb{E}}
\newcommand{\eps}{\varepsilon}
\newcommand{\dd}{\,\mathrm{d}}
\begin{document}
\section*{UnsolvedMath ID: 30004644}
\textbf{Source problem code:} catalogue entry
\subsection{Definitions and mathematical statement}
Let \(T\) be transcendental over \(\Q\), and let \(E/\Q(T)\) be a finite Galois extension with Galois group \(G\). The extension is \(\Q\)-regular if \(E\cap\overline{\Q}=\Q\), inside a common algebraic closure. Its genus is the genus of the smooth projective geometrically integral curve over \(\Q\) with function field \(E\). Assume this genus is \(g\geq 2\). For a finite set \(A\), write \(|A|\) for its cardinality.

The supplied record asks whether, for every such extension,
\[
\lim_{x\to\infty}
\frac{\bigl|\operatorname{Sp}(E)\cap S(G,x)\bigr|}
{\bigl|\operatorname{Sp}(E)^{\mathrm{loc}}\cap S(G,x)\bigr|}
=0.
\]
Here \(\operatorname{Sp}(E)\), \(\operatorname{Sp}(E)^{\mathrm{loc}}\), and \(S(G,x)\) are retained exactly as source notation. The supplied record does not define these sets. In particular, it does not specify the specialization parameters and exclusions, whether specializations must have full Galois group \(G\), the equivalence relation used to count extensions, the local conditions and places involved, or the ordering invariant bounded by \(x\). These choices must be recovered from the cited source before the displayed question becomes a fully specified mathematical statement.

The quotient additionally requires finite numerator and denominator and a positive denominator for the values of \(x\) under consideration. The record provides neither these guarantees nor a convention for zero denominators. No convention is imposed here.

The supplied dated assessment describes a conditional result for abelian groups with sufficiently many branch points and labels the general genus-at-least-two question as remaining open. It does not supply the conditional hypotheses or the required branch-point bound; consequently, that assessment is recorded as a source qualification rather than formulated as a theorem.
\subsection{Short English statement}
For every \(\Q\)-regular Galois extension of \(\Q(T)\) of genus at least two, the source asks whether the displayed ratio of global specialization counts to locally admissible counts tends to zero. The counting sets, local conditions, ordering, and denominator conventions are unspecified in the supplied record, so a precise interpretation still requires source review.
\subsection{Sources}
\begin{itemize}
\item Supplied local UnsolvedMath registry record, permanent ID \(30004644\), title \emph{Local-Global Density of Galois Cover Specializations}, exact source problem code \texttt{catalogue entry}. The record belongs to the Oberwolfach Reports Open Problems collection and supplies proposed year \(2021\); no individual proposer is identified.
\item Supplied report citation: \emph{Homotopic and Geometric Galois Theory (2022)}, Oberwolfach Reports, pp.~663--744, DOI \texttt{10.4171/owr/2021/12}. \url{https://doi.org/10.4171/owr/2021/12}. The background attributes the problem to workshop open-problem discussions documented by Oberwolfach Reports. The title year and DOI are retained as supplied.
\item Supplied literature reference: \emph{Density results for specialization sets}. \url{https://doi.org/10.1017/S1474748019000537}.
\item Supplied literature reference: \emph{Oberwolfach Report 12/2021}. \url{https://doi.org/10.4171/owr/2021/12}.
\item Literature assessment dated exactly \texttt{2026-08-22}: the supplied curation labels the problem partially solved, reports conditional convergence to zero for abelian groups with sufficiently many branch points, and describes the general genus-at-least-two formulation as open. This is supplied dated curation derived from the OWR material, not an independently verified literature assessment. The record expressly requests independent verification of time-sensitive status and cited claims.
\item The registry labels the statement \texttt{corrected\_verified} and reports a manual semantic, typographic, and LaTeX review asserting self-containment and preservation of the target. That supplied review label does not resolve the missing definitions identified in this draft, and has not been independently verified.
\end{itemize}
\end{document}
